\documentclass[11pt,a4paper]{article}
% Strain permits Free Format initial submissions. See SUBMISSION.md.
\usepackage[margin=25mm]{geometry}
\usepackage[T1]{fontenc}
\usepackage{lmodern,microtype,amsmath,amssymb,graphicx,booktabs,array}
\usepackage{xcolor}
\usepackage[numbers,sort&compress]{natbib}
\usepackage[colorlinks=true,allcolors=blue!45!black]{hyperref}
\usepackage{fancyhdr,caption}
\captionsetup{font=small,labelfont=bf}
\pagestyle{fancy}\fancyhf{}\fancyhead[L]{Solver-confirmed initialization for large-deformation correlation}\fancyhead[R]{Draft}\fancyfoot[C]{\thepage}
\setlength{\headheight}{14pt}
\newcommand{\uu}{\mathbf u}\newcommand{\xx}{\mathbf x}\newcommand{\FF}{\mathbf F}
\title{Solver-confirmed initialization and affine propagation for large-deformation digital image correlation on portable graphics hardware}
\author{Author names and affiliations to be confirmed}
\date{Research manuscript draft --- 14 November 2026}
\begin{document}\maketitle
\begin{abstract}
Digital image correlation may fail to converge when two speckle images differ by large displacements and heterogeneous deformation. We investigate initialization separately from measurement: image matching proposes starting states, local photometric solves estimate affine states, and accepted states initialize neighbouring solves. The evaluated chain combines a learned coarse correspondence generator, image-derived orientation and affine hypotheses, local zero-mean normalized cross-correlation, and first-order affine transport. Reliability-guided correlation and feature-assisted affine propagation are established prior art; the contribution examined here is a specified selection/confirmation policy, controlled component ablations and portable implementation evidence. On twenty heterogeneous pairs with 22,416 visible centres, adding propagation to affine plus local-correlation initialization increases correct recovery from 72.56\% to 92.13\%. A separate confirmatory workflow, with neighbour confirmation and consistency checks, recovers 23,570 of 24,641 visible heterogeneous centres without an observed accepted visible displacement error of at least two pixels. These two propagation policies are not interchangeable. Smaller supports recover further centres but worsen gradient error on common correct centres. Severe localized shear remains a failure regime. Real experimental images provide acceptance and remapping evidence without independent accuracy validation. Shared WebGPU kernels support portable matching and solver-related computation, but neither full-pipeline device residency nor subsecond processing is established. The method initializes an existing correlation solver; it does not replace digital image correlation or certify mechanical accuracy from numerical acceptance.
\end{abstract}
\noindent\textbf{Keywords:} digital image correlation; initialization; large deformation; affine propagation; WebGPU.

\section{Introduction}
Digital image correlation (DIC) measures image-plane motion by solving local or global photometric alignment problems. A local solver can be precise near the correct solution and nevertheless fail when its initial state points to unrelated material. This is particularly important for two widely separated frames without temporal information. Large uniform translation demands correspondence discovery; strong variation of the displacement gradient also challenges the validity of a local affine model. These are different obstacles.

The practical question is therefore not whether a matcher independently measures every requested point. It is whether it supplies enough useful starting states for the correlation solver to recover the material regions of interest. Our organizing principle is: find proposals, solve a few windows, then transport their affine states to initialize further photometric solves. A proposal or interpolated state is not a DIC measurement. A numerically accepted solve is not automatically physically correct.

This principle has substantial DIC precedent. Pan's reliability-guided method prioritizes propagation through computed points with high zero-mean normalized cross-correlation \citep{pan}. Zhou, Pan and Chen combine feature correspondences, affine seed initialization and transfer of refined parameters to adjacent points \citep{zhou}. Their publisher abstracts establish these mechanisms; the present comparison does not infer unreported implementation details from those abstracts. Dense DAISY description with structured search has also been investigated for point-wise DIC initialization \citep{li}. We do not claim novelty for the succession matching--affine initialization--propagation.

The contribution assessed here is narrower: an explicit image-derived multi-hypothesis selection and neighbour-confirmation policy, its separation from a simpler historical propagation baseline, a controlled ablation of local correlation and a learned residual field, and numerical preservation in portable graphics implementations. The article first specifies the retained chain, then evaluates recovery and failure regimes. Explored descriptors and interpolation alternatives are separated in the supplement, rather than presented as consecutive modules. Strong-strain support mismatch explains several negative alternatives, but does not establish universal superiority of one matcher.

The principal evidence is synthetic, with an independent confirmatory cohort and unchanged evaluation thresholds. The experimental pair is shown to make the procedure tangible, not to substitute visual agreement for displacement truth. All numerical claims retain their campaign denominator and implementation scope.

\section{Kinematics, photometry and the initialization problem}
Let $I_0$ and $I_1$ be monochrome reference and deformed images. The material image transformation is $\boldsymbol\varphi(\mathbf x)=\mathbf x+\mathbf u(\mathbf x)$, with gradient $\mathbf F=I+\nabla\mathbf u$. At requested centre $\mathbf x_i$, the local affine approximation is
\begin{equation}
\boldsymbol\varphi_i(\boldsymbol\xi)=\mathbf x_i+\mathbf d_i+(I+\mathbf G_i)\boldsymbol\xi,
\qquad \mathbf p_i=[u,u_x,u_y,v,v_x,v_y]^T,
\label{eq:warp}
\end{equation}
where $\boldsymbol\xi$ is position relative to the reference centre, $\mathbf d_i=(u,v)^T$ is its displacement and $\mathbf G_i=\left[\begin{smallmatrix}u_x&u_y\\v_x&v_y\end{smallmatrix}\right]$ its displacement gradient. Image coordinates and displacements are in pixels; the gradients are dimensionless when both images use the same pixel unit. These are image-plane kinematics. Their mechanical interpretation requires appropriate calibration and imaging assumptions. In particular, $\mathbf G$ is not itself a finite-strain tensor under large rotations.

For a smooth transformation, $\boldsymbol\varphi(\mathbf x_i+\boldsymbol\xi)=\boldsymbol\varphi_i(\boldsymbol\xi)+O(\|\boldsymbol\xi\|^2)$. A reference square and a target square at the correct centre need not contain the same material. Under 60\% horizontal extension, a 64-pixel material support occupies about 102 pixels horizontally in the target. Scale and orientation normalization do not generally compensate anisotropic stretching and shear. Geometrically adapted descriptors exist; the difficulty concerns comparison without the appropriate support compensation, not an absolute inability of descriptors to handle deformation.

\begin{figure}[tb]\centering\includegraphics[width=\linewidth]{figures/material_support.pdf}
\caption{Pedagogical support geometry, not experimental data. A reference square maps to a parallelogram under the prescribed affine transformation. A fixed target square at the correct centre samples different material. Sampling the target on the transformed support and expressing it in reference coordinates restores comparable material supports. The centre alone does not determine this correction.}\label{fig:support}\end{figure}

Under ideal conserved brightness and differentiability,
\begin{equation}
I_0(\mathbf x)=I_1(\boldsymbol\varphi(\mathbf x)),\qquad
\nabla I_0(\mathbf x)=\mathbf F(\mathbf x)^T\nabla I_1(\boldsymbol\varphi(\mathbf x)).
\label{eq:advection}
\end{equation}
Filtering, interpolation and illumination changes make this relation approximate in real images. Its purpose here is to explain support and gradient mismatch, not to introduce a new loss into the retained method.

For sampled reference intensities $\mathbf f$ and warped target intensities $\mathbf g(\mathbf p)$, subtract the respective means and normalize each centred vector to unit norm. With non-zero norms, denoting these vectors by hats gives the zero-mean normalized sum of squared differences (ZNSSD):
\begin{equation}
E_i(\mathbf p)=\|\widehat{\mathbf f}-\widehat{\mathbf g}(\mathbf p)\|^2
=2\{1-\widehat{\mathbf f}^T\widehat{\mathbf g}(\mathbf p)\},\qquad
c_i=\widehat{\mathbf f}^T\widehat{\mathbf g}(\mathbf p).
\label{eq:znssd}
\end{equation}
Thus minimizing ZNSSD and maximizing zero-mean normalized cross-correlation (ZNCC) are equivalent only for identical valid samples and normalizations. A constant subset does not define this ideal normalization.

The inverse-compositional Gauss--Newton (IC-GN) method linearizes the reference-side incremental warp and composes its inverse with the current warp \citep{baker,panfast}. The six unknowns in equation~\eqref{eq:warp} are coupled through the texture gradients. A well-conditioned local problem can still have a narrow convergence basin: an unrelated target patch may lead to another stationary point. Initialization is the search for a starting state in the useful basin of the fixed solver. It does not replace that solver. The supplementary derivation explains the normalized differential; the actual frozen solver settings are reported separately.

We distinguish three outcomes throughout: numerical acceptance by the solver/policy, displacement correctness relative to independent analytical truth, and metrological accuracy of displacement and gradient. The two-pixel scoring threshold used below is a qualification convention for recovery, not a measurement uncertainty.

\section{From two images to solver-accepted states}
\label{sec:method}
\subsection{Objects and scope of the retained workflow}
A \emph{requested centre} is fixed by the DIC sampling grid. A \emph{correspondence candidate} is a pair of coordinates proposed by image matching. A fitted affine map predicts a position; it does not observe that position independently. A local correlation maximum is a \emph{photometric proposal}; its six-parameter conversion is an \emph{initial state}. A \emph{solver-accepted seed} has completed a local solve with acceptable numerical status. In the confirmatory workflow, an \emph{anchor} additionally passes neighbour confirmation. None of these definitions implies independently established correctness.

Figure~\ref{fig:pipeline} describes the confirmatory workflow. The historical causal ablation uses the same proposal family but a simpler propagation rule, explicitly distinguished below. Analytical displacement and visibility are excluded from operational selection. They enter training/calibration on their own data, or scoring after inference on test data. A supplied material mask, unlike analytical visibility, may be an operational input.

\begin{figure}[tb]\centering\includegraphics[width=\linewidth]{figures/pipeline.pdf}
\caption{Confirmatory operational chain. Boxes distinguish proposals from solved states. Reduced-raster matching and local correlation initialize native-pixel solves. The feedback loop transports accepted states only to initialize new solves. Analytical truth belongs to the separate evaluation block. Conditional regional retries use image-derived hypotheses, never analytical error.}\label{fig:pipeline}\end{figure}

\subsection{Learned coarse correspondence generator}
The main campaigns use the project's \texttt{NDFPairMatcher}. Here NDF names the research neural deformation field model. Its shared convolutional encoder describes the frames independently; it is not a cross-attention Transformer. Each frame is centred and divided by its root-mean-square intensity variation (floor $10^{-4}$). Two stride-two convolutions and a final convolution form a stride-four feature map. The tested contextual option adds spatial convolutions, including a dilated layer, \emph{within each image}. Each output has 32 components normalized to unit Euclidean norm.

For the $192\times108$ input, there are $48\times27=1296$ feature locations. With source feature $\mathbf f_i$ and target feature $\mathbf h_j$, the model calculates
\begin{equation}
s_{ij}=\mathbf f_i^T\mathbf h_j,\quad
\pi_{ij}=\frac{\exp(s_{ij}/T)}{\sum_k\exp(s_{ik}/T)},\quad
j_i=\arg\max_j\pi_{ij}.\label{eq:coarse}
\end{equation}
The learned temperature is clamped to $[0.025,0.5]$. Matching is global on this reduced grid. Hard target coordinates feed the affine fit. The model also returns confidence $\max_j\pi_{ij}$ and a soft expected position. A separate residual head combines source features, probability-weighted target features, soft displacement and confidence. It is used in the residual ablation, not in the retained zero-residual field. Omitting the NDF residual does not remove the trained encoder or temperature.

The checkpoint comes from the local 60-epoch experiment with 128 training pairs, seed 20261114, contextual features and physical-loss weight 0.02. Its 32 validation pairs use distinct texture identities according to the training manifest. The checkpoint path, present-file fingerprint and training protocol accompany the supplement. Validation accuracy of this network alone is not the performance of the complete initializer.

\subsection{Orientation hypotheses and affine preconditioning}
The archived reference bank has ten angles: $-180,-140,-100,-60,-20,0,20,60,100,140$ degrees. It is a nine-angle uniform circle with identity inserted, not a ten-angle uniform circle. For centre $\mathbf c=((W-1)/2,(H-1)/2)$ in reduced-raster pixels, the tested target is
\begin{equation}
I_{1,\alpha}(\mathbf z)=I_1\bigl(\mathbf c+R_\alpha(\mathbf z-\mathbf c)\bigr).\label{eq:rotation}
\end{equation}
Coordinates increase rightward and downward; the equation fixes the sign. Bilinear sampling fills missing weight with intensity 0.5. Matched target locations are mapped through the same coordinate map before fitting source-to-original-target geometry.

An image-derived probe uses a spatial grid with budget 256 queries (252 actual queries at this aspect ratio) against all 1296 target candidates. Mutual maxima are checked within this probe matrix. Global affine consensus ranks each angle by supported-query fraction times mean supported confidence; the fraction uses the nominal budget 256. It makes 32 random triplet trials, rejects determinant magnitudes below $10^{-7}$, uses a five-pixel residual threshold in the reduced raster, and refits when at least six points support the best model. A non-identity winner requires at least 24 supported queries and a relative score gain of 0.1 over identity. The diagnostic true angle is never a selection input.

For the selected orientation, all 1296 hard matches are mapped back and fitted by a separate global affine estimator: 128 triplet trials, a five-pixel threshold, at least 24 inliers, then eight Huber-type reweighted least-squares steps. The fitted map $A\mathbf z+\mathbf t$ must cover four cells of a $3\times3$ source partition, have $0.05<\det A<20$, smallest singular value above $10^{-6}$ and condition number below 20. The historical ablation also rejects $\|A-I\|_F<0.03$. The later confirmatory campaign explicitly retains near-identity maps when the residual is disabled, retaining pure translations. This change is not attributed retroactively to the ablation.

An accepted affine map defines the rectified target $J_1(\mathbf z)=I_1(A\mathbf z+\mathbf t)$. If rejected, the selected rotated target remains and the residual starts at zero in that frame. The inverse-compositional Gauss--Newton (IC-GN) solver subsequently sees the original native images, not these rectified images.

\subsection{Local photometric proposals: translation at fixed geometry}
The archive labels this component NCC (normalized cross-correlation); its cost is \emph{zero-mean} normalized cross-correlation (ZNCC). At each stride-four location $\mathbf z$, the no-residual baseline compares a $9\times9$ reference support with target supports centred at $\mathbf z+\boldsymbol\delta$ in the rectified raster. Each component of $\boldsymbol\delta$ ranges from $-8$ to $8$ integer pixels. For valid samples,
\begin{equation}
c(\boldsymbol\delta)=\frac{\sum_k(f_k-\bar f)(g_k(\boldsymbol\delta)-\bar g)}
{\sqrt{\sum_k(f_k-\bar f)^2\sum_k(g_k(\boldsymbol\delta)-\bar g)^2}},\quad
\boldsymbol\delta_* = \arg\max_{\boldsymbol\delta\in\{-8,\ldots,8\}^2}c(\boldsymbol\delta).
\label{eq:localncc}
\end{equation}
Only translation is searched. Local geometry is identity \emph{in the globally rectified raster}; there is no joint optimization of six affine parameters. Bilinear sampling is used; candidates extending beyond the raster receive score $-2$. The shader floors variances and denominator. A constant patch therefore has a numerical score rather than a trustworthy correspondence: there is no independent contrast-confidence gate here. A deterministic maximum is selected, followed by a parabolic correction along each axis, limited to half a pixel when the maximum is interior. Scores, margins and validity are returned, but the qualifier uses the displacement without an additional score or margin rejection. The native solver supplies acceptance.

The residual is recomposed in the original target frame:
\begin{equation}
\widehat{\boldsymbol\varphi}(\mathbf z)=A[\mathbf z+\boldsymbol\delta_*(\mathbf z)]+\mathbf t.
\label{eq:composition}
\end{equation}
A bilinear field evaluates displacement and derivatives at native centres after half-pixel-aware conversion. At isotropic scale 20, one reduced residual pixel corresponds to 20 native pixels; this is a unit conversion, not an accuracy claim. Appendix~\ref{sec:coordinates} gives the anisotropic formula. Neither interpolation nor a correlation maximum is reported as a new DIC measurement.

\subsection{Native solve, hypothesis selection and anchor confirmation}
Confirmatory settings are a 64-pixel square support, 64-pixel centre spacing, affine motion, linear interpolation, at most 30 iterations, tolerance $10^{-3}$ and minimum solver ZNCC 0.65. The optional outlier policy is disabled. States with status greater than 0.5 and finite parameters are confirmation candidates. Current solver inspection shows convergence judged by incremental corner change, then correlation threshold; contrast, singular updates and invalid supports can also cause rejection. The archives do not fingerprint the complete external solver build. Recorded settings are verified; the exact historical low-level rejection sequence cannot be asserted to equal today's external source.

The confirmatory hypothesis list is ordered: zero displacement, three separately generated ranked-rotation hypotheses, global affine, and affine plus local correlation. Each is solved on every second row and column of the requested grid, restricted to the operational mask. The winner has the most finite solver-accepted states; ties retain the earlier list entry. This native sparse test is distinct from coarse orientation ranking and has no fixed 128-window cap in this campaign.

Confirmation visits selected sparse candidates in deterministic order. A candidate and an already verified neighbour are confirmed when affine transport predicts the neighbour within 8 native pixels in displacement and 0.25 in gradient norm. Still-unverified valid neighbours are also solved from transported states: at least two consistent accepted solves confirm the parent and those children. Thus an existing consistent pair can be confirmed without two newly solved neighbours. The variable ``confirmed seeds'' includes children solved in this stage; it is not an initial match count. Thresholds were selected on development data and frozen before confirmation.

\subsection{Two archived propagation policies}\label{sec:policies}
For accepted parent $j$ and neighbouring centre $i$, transport supplies
\begin{equation}
\mathbf d^{0}_{i\leftarrow j}=\mathbf d_j+\mathbf G_j(\mathbf x_i-\mathbf x_j),\qquad
\mathbf G^{0}_{i\leftarrow j}=\mathbf G_j.\label{eq:transport}
\end{equation}
The neighbour is then solved on the images. For a pedagogical example, $\mathbf G_j=\operatorname{diag}(0.1,0)$ and a 64-pixel horizontal separation add 6.4 pixels to horizontal displacement. Copying translation misses this increment even if the original displacement is hundreds of pixels. This is an analytical illustration, not a benchmark.

In the \textbf{historical causal ablation}, every one-pass solver-accepted state seeds a synchronous frontier. Each unresolved centre receives the componentwise median of states transported from the accepted eight-neighbour snapshot. It is attempted at most three times. Stop occurs when the eligible frontier is empty, not necessarily at the first pass with no acceptance. Accepted states are immutable; there is no extra consistency threshold.

In the \textbf{confirmatory workflow}, confirmed anchors seed the frontier. For each unresolved valid centre, the first two accepted neighbours in stable order (up, left, right, down, then diagonals) are available from a fixed pass snapshot. The first is tried, then the second only if that centre was not added. A finite solver-accepted result must remain within 8 pixels and 0.25 gradient norm of the proposal. New states become sources next pass. Stop occurs on an empty frontier or a pass with no additions. There is \emph{no three-attempt lifetime cap} and no median aggregation. Accepted states remain immutable. Algorithm~\ref{alg:frontier} specifies this retained policy.

For a smooth true field, Taylor expansion gives
\begin{equation}
\mathbf u(\mathbf x_j+\boldsymbol\eta)-[\mathbf d_j+\mathbf G_j\boldsymbol\eta]
=\mathbf e_j+E_j\boldsymbol\eta+O(\|\nabla^2\mathbf u\|\|\boldsymbol\eta\|^2),\label{eq:taylor}
\end{equation}
where $\mathbf e_j$ and $E_j$ are the parent's displacement and gradient errors. Transport removes the leading variation missed by translation copying, but not seed error, curvature or jumps. It is neither final interpolation nor a mechanical-equilibrium solution. It differs from PatchMatch candidate propagation and temporal prediction.

\begin{table}[p]\small\centering
\caption{Complete confirmatory workflow. Every solve uses native images and unchanged local solver settings.}\label{alg:chain}
\begin{minipage}{.95\linewidth}\begin{enumerate}
\item Fix centres, support and supplied mask. Reduce images; extract coarse features.
\item Rank ten image-derived orientations. Fit global affine geometry; construct affine and local-correlation states. Load three separately ranked rotation proposals.
\item Solve each hypothesis on the same sparse valid grid. Select largest finite accepted count, breaking ties by list order.
\item Confirm anchors by neighbour consistency. Run Algorithm~\ref{alg:frontier}; record parents, generations and attempts.
\item If an operational-mask component has fewer than three anchors, retry unresolved components of at least 16 centres using the same hypothesis list. Merge only new acceptances; stop retries if no progress.
\item If the minimum-anchor condition still fails, fit local affine hypotheses from archived hard matches around unresolved centres: 192-pixel radius, four-pixel residual threshold, six inliers, eight candidates, inlier fraction 0.15, 256 trials. Repeat selection/confirmation, preserving accepted states. This branch is conditional.
\item Return accepted states and unresolved/rejected states separately. Only now score synthetic correctness and analytical visibility.
\end{enumerate}\end{minipage}\end{table}

\begin{table}[p]\small\centering
\caption{Snapshot frontier in the confirmatory workflow, corresponding to \texttt{seed\_propagation.propagate}. The historical median/three-attempt rule is separate.}\label{alg:frontier}
\begin{minipage}{.95\linewidth}\begin{enumerate}
\item Set $S$ to confirmed anchors; initialize attempts and provenance.
\item Copy $S$ to source snapshot $S_0$. For each valid $i\notin S_0$, select at most two neighbours in $S_0$ in fixed order. Stop if no centre has a source.
\item For source rank 1 then 2, batch centres not added earlier this pass. Transport sources by equation~\eqref{eq:transport}; increment attempted-centre counters; call the solver.
\item Add a state only if solver status is accepted, parameters finite, and displacement/gradient changes from its proposal do not exceed the frozen limits. Record parent and pass. Never replace a previously accepted state.
\item Stop if this pass added no state; otherwise repeat step 2. Return unresolved states and rejection reasons, never treating a filled prediction as a measurement.
\end{enumerate}\end{minipage}\end{table}

\subsection{Operational masks are not analytical visibility}
The real-image mask restricts solves through the solver's region-of-interest selection. Without a mask all requested centres enter initialization; synthetic visibility is used only for scoring. Host-side neighbours check grid adjacency and valid endpoint states, not the pixel segment between centres or the solver's prepared region links. Diagonals can bridge narrow masked separations. Fracture-safe connectivity is therefore not claimed for these archives. Analytical interface labels in the support study are post-hoc diagnostics, not an operational crack detector. Later solver changes do not retroactively alter these measurements.

\section{Implementation and reproducibility}
The operational objects have different resolutions. In the main synthetic campaigns both native images are $3840\times2160$. Area reduction produces $192\times108$ rasters; the convolutional feature grid is $48\times27$. Rotation ranking, global affine fitting and the $9\times9$ local photometric supports operate in this reduced raster. Native IC-GN supports are $64\times64$ with 64-pixel centre spacing. The later 32-pixel support experiment preserves these centres; it is not an increase in matching resolution or a denser sampling grid.

The core proposal/solve interfaces are summarized in Table~\ref{tab:interfaces}. The supplementary parameter tables and machine-readable campaign configurations retain exact recorded values. Parameters absent from an archive are explicitly left unknown rather than copied from another experiment. The checkpoint fingerprint identifies the file inspected for this revision; it cannot prove which unrecorded historical binary was loaded.

\begin{table}[tb]\centering\small
\caption{Interfaces in the confirmatory chain. Pixels refer to the stated raster; all rows are image-derived until the separate scoring stage.}\label{tab:interfaces}
\begin{tabular}{p{.19\linewidth}p{.34\linewidth}p{.36\linewidth}}\toprule
Stage & Input and operation & Output / selection\\\midrule
Coarse discovery & $192\times108$ images, shared encoder, global feature similarities & $1296$ hard coordinate pairs; confidence\\
Orientation and geometry & Ten angles, sparse reciprocal probes, affine consensus & Image-derived ranking and global map\\
Local proposal & Rectified raster, $9\times9$ ZNCC, residual radius 8 & Two-component residual; affine composition\\
Grid conversion & Bilinear field and analytic cell derivatives & Six native parameters at requested centres\\
Selection and confirmation & Native sparse solves, neighbour re-solves & Finite accepted states, confirmed anchors\\
Propagation & Native pixels and transported affine states & Re-solved accepted field, parents, generations, rejects\\\bottomrule
\end{tabular}\end{table}

WebGPU is a portable graphics-compute interface, and WGSL denotes its WebGPU Shading Language \citep{webgpu,wgsl}. Canonical shaders in the matching package are shared by Python and browser-facing code; they do not imply identical host orchestration. Python supplies control, model inference and data preparation. The archived coarse learned matcher uses PyTorch on the central processing unit (CPU); the local-correlation and solver kernels use the graphics processing unit (GPU). Global affine fitting and the original frontier involve host arrays. Later device-resident frontier experiments keep local states and active masks on one device, with compact counters read per pass, but do not establish residency of the entire learned-matching chain. An Open Neural Network Exchange (ONNX) Runtime WebGPU prototype is a separate execution experiment, not the backend silently attributed to the main coverage tables.

This revision freezes the inspected working-tree sources and their SHA-256 cryptographic fingerprints, the base Git commit, configurations and source-result paths. The repository is based on \texttt{4dce5d0} with later uncommitted changes; that commit alone is not a reproducible release of these studies. Current code establishes implemented behaviour; archived configurations establish requested campaign settings; saved window arrays establish retained outputs. Historical equivalence between those levels remains incomplete where source or external-solver fingerprints were not recorded. No algorithm, threshold, checkpoint or archived output was altered for this editorial revision.

\section{Evaluation protocol}
The experiments answer distinct questions, summarized in Table~\ref{tab:campaigns}. The principal comparison is not a uniform tournament of every explored matcher. Calibration uses analytical labels on development data to freeze selection/consistency settings. Test inference receives images and, when supplied, a real material mask; analytical visibility and deformation enter only subsequent scoring. Existing historical cases are regression diagnostics, not newly independent tests.

\begin{table}[tb]\centering\small
\caption{Campaign identities and populations. All synthetic native rasters are $3840\times2160$. The complete protocol dictionaries, identifiers and fingerprints are supplied separately.}\label{tab:campaigns}
\begin{tabular}{p{.22\linewidth}p{.24\linewidth}p{.42\linewidth}}\toprule
Campaign & Population & Question / key distinction\\\midrule
Causal ablation & 20 heterogeneous pairs; groups 16000000, 17000000 & Four residual/correlation variants; status-only median propagation; 22,416 visible centres\\
Confirmation & 20 heterogeneous and 20 translation pairs; groups 23000000, 24000000 & Frozen sparse confirmation and consistency-gated propagation; near-identity affine retention enabled\\
Support transfer & Five pairs in each of five families; group 62000000 & 64 to 32 pixels, identical centres; derivative/coverage tradeoff\\
Severe shear & Three textures at each peak gradient 0.33, 0.67, 2.0; group 59000000 & Exact-centre analytical scoring; fixed 64-pixel support failure\\
Experimental & Adrien frames 000000/000075; supplied mask & $7920\times6004$ images; acceptance and remapping only\\
Implementation & Four-pair GPU parity and isolated cached-script comparison & Numerical preservation and scoped timing, not new accuracy cohorts\\\bottomrule
\end{tabular}\end{table}

Let $V_i$ indicate analytical centre visibility, $S_i$ policy acceptance, and $e_i=\|\mathbf d_i-\mathbf u_{\rm true}(\mathbf x_i)\|$. Correct recovery uses the \emph{strict} threshold $e_i<2$ native pixels:
\begin{equation}
C=\frac{\sum_i V_iS_i\mathbf1[e_i<2]}{\sum_i V_i},\quad
N_{\rm false}=\sum_i V_iS_i\mathbf1[e_i\geq2],\quad
Q=\frac{\sum_i V_iS_i\mathbf1[e_i<2]}{T_{\rm complete}}.
\label{eq:metrics}
\end{equation}
Non-finite accepted states, rejected states and accepted invisible centres are separately audited. $T_{\rm complete}$ must include input preparation, hypotheses, failed trials, propagation and result recovery. Where that time is absent, $Q$ is not reported as end-to-end throughput. Descriptor diagnostics at five pixels are not pooled with these DIC outcomes. An analytical local-affine solve is a diagnostic reference, not an absolute mathematical upper bound.

Reanalysis for this revision reads saved arrays, recomputes errors from states and truth where both exist, and checks the reported counts. It is not a new independent experiment. All four causal denominators remain 22,416. Confirmation has 24,641 visible heterogeneous and 29,459 visible translation centres. Pair-level mean gains and pair-bootstrap intervals are distinct from ratios of pooled counts. The support study has only five pairs per family and can reuse texture identities between families, so its cross-family samples are not independent. Its gradient errors are compared on centres correct at both support sizes; newly recovered centres are counted separately.

Synthetic families preserve material texture under the generator's deformation. Exact-centre analytical truth is used in the corrected severe-shear and support studies; historical arrays are not silently regenerated with today's generator. The discontinuous-opening support examples are grid aligned and cannot establish arbitrary crack robustness. Real-image observations have no independent displacement truth. The experimental acceptance denominator is the supplied mask, not a synthetic visibility mask.

\section{Results}
\subsection{Which component produces the improvement?}
The most informative comparison holds the coarse matcher, orientation bank, affine prefit, twenty heterogeneous pairs and native IC-GN settings fixed. Only the nonlinear residual and local correlation are changed. The historical status-only propagation policy of Section~\ref{sec:policies} is used here. Table~\ref{tab:causal} reports the original one-pass results and a subsequent propagation experiment using those accepted states. All rows share 22,416 visible requested windows. The coarse matcher, prefit and NCC operate on a $192\times108$ raster; the correlation solver uses native $3840\times2160$ pixels. Thus ``affine only'' does not mean ``without a learned matcher''.

\begin{table}[tb]\centering\small
\caption{Controlled component ablation before and after status-seeded propagation. All rows share the same image pairs and visible denominator. Analytical local-affine initialization yields 21,367 correct windows. A seeded pair contains at least one solver-accepted initial state; it need not contain a seed in every disconnected material region.}\label{tab:causal}
\begin{tabular}{lrrrr}\toprule
Initialization & One pass & Propagated & Seeded pairs & False accepted\\\midrule
Affine & 2,070 & 19,915 &18/20&1\\
Affine + NCC &16,266&20,651&19/20&0\\
Affine + NDF &1,748&19,070&17/20&0\\
Affine + NDF + NCC &13,172&19,917&18/20&0\\\bottomrule
\end{tabular}\end{table}

Adding NCC to affine initialization recovers a net 14,196 correct windows in one pass: pooled coverage rises from 9.23\% to 72.56\%. After propagation, however, the corresponding gap is only 736 windows, or 3.28 percentage points (88.84\% versus 92.13\%). This difference changes the interpretation of the component. NCC is valuable not only for immediate coverage, but also for providing an additional seeded pair and reducing propagation effort: median propagation rings fall from 13 to 3 in this implementation. The baseline must therefore include propagation before claiming that a more expensive initializer is necessary.

Propagation adds 4,385 correct windows to affine + NCC, a 19.56-point pooled gain, without a new visible false acceptance. Affine alone gains 17,845 windows, but retains one false initial acceptance. That state has a 9.466-pixel error on pair 16051009/16091009 and is already present in the initial pass. Its presence is reported rather than removed using truth. No accepted invisible windows occur in these four propagation comparisons.

Adding the tested NDF residual when NCC is present loses 3,094 one-pass successes. Windowwise accounting identifies 245 successes unique to the residual and 3,339 successes lost to it. The mean paired coverage change is $-13.15$ points, with a 95\% pair-bootstrap interval of $[-16.43,-9.69]$. After propagation the residual still yields 734 fewer correct windows and one fewer seeded pair. The operational baseline therefore omits this residual. This conclusion does not reject an untested network explicitly trained for residual prediction after rectification.

\begin{figure}[tb]\centering\includegraphics[width=\linewidth]{figures/causal_components.pdf}\caption{Matched component ablation: one-pass and propagated correct coverage on the same 22,416 visible windows. The dashed line is analytical local-affine initialization, a diagnostic reference. Neural deformation field (NDF) denotes the tested residual checkpoint. The large effect of propagation is distinct from the smaller remaining difference between initializers.}\label{fig:causal}\end{figure}


\subsection{Held-out heterogeneous and translation fields}
The confirmatory heterogeneous campaign recovered 23,570 of 24,641 visible windows (95.65\%), with no false accepted visible windows. The paired mean gain over the archived baseline was 23.65 percentage points, with a 95\% pair-bootstrap interval of 14.78--35.19 points. Against a matched re-solve control, the mean gain was 23.40 points (14.69--34.61). Figure~\ref{fig:paired} shows why the pooled count is insufficient: difficult pairs account for substantial recovery, and a small aggregate difference can conceal a local failure.

The translation campaign recovered 28,641 of 29,459 visible windows, compared with 25,173 for the corrected baseline, with no pair regression and no visible false acceptance. An earlier comparison had rejected near-identity affine prefits and was not used as the confirmatory baseline. Retaining that correction is essential because pure translation is a basic control for the geometric preprocessing itself.

\begin{figure}[tb]\centering\includegraphics[width=.92\linewidth]{figures/paired_coverage.pdf}\caption{Paired heterogeneous-field coverage reconstructed from the archived confirmatory report. Lines connect the baseline and propagated result on each image pair, sorted by gain. Both methods use that pair's visible requested centres as denominator. These are final correlation outcomes, not direct feature-match counts.}\label{fig:paired}\end{figure}

\subsection{Independent same-centre transfer}
The five-family experiment changes support size while preserving centres. Table~\ref{tab:multi} and Figure~\ref{fig:multi} show increased coverage in every pooled family, with no loss of previously correct centres and no visible false acceptance in these draws. Nevertheless, gradient error increased on the common correct population in every family. In heterogeneous fields, displacement root-mean-square error decreased from 0.0589 to 0.0191 pixels, while gradient root-mean-square error increased from 0.0009 to 0.0025. Smaller support can reduce non-affine mismatch while containing less information for derivative estimation.

\begin{table}[tb]\centering\caption{Independent transfer at identical centres, five pairs per family. Correct means an accepted solve with displacement error below two pixels.}\label{tab:multi}
\begin{tabular}{lrrr}\toprule Family & Visible & 64-pixel correct & 32-pixel correct\\\midrule
Translation&6467&6322&6395\\Rigid rotation&6278&6008&6144\\Affine&5845&5616&5718\\Heterogeneous&6061&5804&5920\\Discontinuous&6556&6414&6443\\\bottomrule\end{tabular}\end{table}
\begin{figure}[tb]\centering\includegraphics[width=\linewidth]{figures/multifamily.pdf}\caption{Coverage after smaller-support transfer at the same requested centres. Bars start at zero to preserve the scale of improvement. No visible false acceptance occurs in these draws; this does not establish universal safety. The severe-shear qualification remains failed.}\label{fig:multi}\end{figure}

Figures~\ref{fig:spatial} and~\ref{fig:discontinuous} connect the window-level statistics to the actual speckle images. Each shows the full image pair, the spatial outcome, and native-pixel supports before and after the archived local affine correction. The local example is the valid correct centre nearest the image centre, excluding interface-crossing supports; it is not selected by the smallest photometric residual. The pullback uses the estimated six parameters, never the analytical displacement. Full-frame views are reduced for printing, whereas the enlarged supports retain native sampling. The absolute intensity difference is a visual diagnostic, not the normalized acceptance criterion or a displacement-error map. No gain or offset is fitted for these displays; illumination changes can therefore leave a non-zero intensity difference despite successful geometric registration.

\begin{figure}[p]\centering\includegraphics[width=\linewidth]{figures/spatial_transfer.pdf}
\caption{Synthetic heterogeneous pair 62051005/62091005, at $3840\times2160$ pixels. (a,b) Actual reference and deformed speckle images. (c) Archived outcome at the requested centres: blue is correct at both support sizes, green is newly recovered at 32 pixels, and orange is unresolved or incorrect among visible centres. (d) A native 32-pixel reference support, located by the magenta cross. (e) The target support centred at the estimated corresponding position, with translation correction only. (f) The target pulled back using the complete estimated local affine transformation. (g) Absolute intensity difference between (d) and (f), on a fixed 0--0.25 scale; larger differences saturate. The same grayscale range is used for (d)--(f). No image is warped with analytical truth. These are final solver outputs after initialization and transfer, not initializer-only precision.}\label{fig:spatial}\end{figure}

\begin{figure}[p]\centering\includegraphics[width=\linewidth]{figures/discontinuous_speckle.pdf}
\caption{Synthetic discontinuous pair 62051004/62091004. Panels, scales and outcome definitions match Figure~\ref{fig:spatial}. The full images show the separated material regions; no smooth global remapping is used to bridge the interface. The enlarged 32-pixel support belongs to one continuous region and therefore does not demonstrate successful correlation across a crack. This archived family has a grid-aligned opening, a limitation of the experiment.}\label{fig:discontinuous}\end{figure}
\subsection{Recorded stages and propagation provenance}
Figure~\ref{fig:chainexample} follows one existing heterogeneous pair through the recorded proposal and solve states. The pre-correlation affine state and the post-correlation initial state are archived separately. The plot does not reconstruct an unavailable raw score volume or claim that a projected centre was directly matched. The illustration is a secondary archive analysis, not an additional validation pair.

\begin{figure}[p]\centering\includegraphics[width=\linewidth]{figures/recorded_chain.pdf}
\caption{Recorded stages for heterogeneous pair 62051005/62091005. (a) Thirty-two evenly indexed hard candidates from the archived 1296 pairs, shown as arrows on the native reference; arrows are proposals, not verified matches. (b) Native 64-pixel reference support at the solver-accepted centre nearest the image centre. (c) Target pullback with archived affine-only initial parameters. (d) Target pullback with archived affine-plus-local-correlation initial parameters evaluated at that centre; this is the interpolated proposal, not a reconstructed raw correlation maximum. (e) Pullback after the archived native IC-GN solve. (f) Absolute final intensity difference. All transformations use saved estimated states, never truth. This is the 64-pixel initial solve, distinct from the 32-pixel support-transfer illustration.}\label{fig:chainexample}\end{figure}

\begin{figure}[p]\centering\includegraphics[width=\linewidth]{figures/propagation_limits.pdf}
\caption{Secondary archive analysis of propagation and its limits. Top: heterogeneous pair 24051003/24091003, with nine confirmed anchors. Bottom: shear pair 59051005/59091005 at peak gradient 2.0. Left: confirmed-anchor locations; middle: recorded accepted generations; right: final analytical scoring (strict error below two native pixels), with rejected visible centres and false acceptances distinguished by symbols. The severe-shear row contains 22 false acceptances; numerical acceptance did not establish correctness. Exact-centre truth is used only for this scoring panel. These selected cases illustrate mechanisms and are not estimates of family-wide performance.}\label{fig:propagation}\end{figure}

\subsection{Failure envelope and experimental images}
At peak shear gradient 2.0, the frozen severe-shear analysis recorded 57 incorrect visible acceptances among 4,080 accepted windows (1.40\%). All 57 were also incorrect with analytical local-affine initialization. Thus the local model, support or solver limits performance as well as the initializer. Better seed positions alone cannot be assumed to remove this failure. Numerical acceptance is not a guarantee of physical correctness.

A separately frozen image-residual suspicion rule detected 119 of 120 false accepted coarse states on an additional heterogeneous and extreme-shear sample, but also flagged 221 of 6,979 correct states. Guarded smaller-support retries recovered 7,070 correct and one false result among 8,075 visible shear centres, leaving 199 previously correct flagged centres unresolved. This is a tradeoff, not a universal error detector. An unknown discontinuity is especially difficult because a nominal neighbourhood can combine two material motions.

Figure~\ref{fig:experimental} shows the experimental pair supplied by Adrien (frames 000000 and 000075), recorded at $7920\times6004$ pixels. Of 7,280 requested centres valid under the archived region-of-interest mask, 7,097 were accepted by the solver (97.49\%). This denominator excludes masked-out centres from the 11,439-centre rectangular grid. The enlargement is the accepted centre nearest the image centre, selected without ranking photometric residuals. The 64-pixel pullback uses its archived final affine parameters. This real-image example uses the masked propagation campaign, whereas Figures~\ref{fig:spatial}--\ref{fig:discontinuous} use the separate 32-pixel transfer campaign; they are not a matched accuracy comparison.

\begin{figure}[p]\centering\includegraphics[width=\linewidth]{figures/experimental_speckle.pdf}
\caption{Experimental images and local remapping, Adrien frames 000000--000075. (a,b) Native-image fields of view, reduced for print. (c) Green marks solver-accepted centres inside the supplied mask; orange marks masked-domain centres not accepted. Unmarked regions are not assessed by this map. There is no independent displacement truth. (d--f) A 64-pixel reference support, target support with translation correction only, and target support with the complete estimated local affine correction. (g) Absolute intensity difference, with the same fixed scale as the synthetic examples. Agreement of speckle patterns demonstrates visual registration at this location, not calibrated displacement or strain accuracy. No global field interpolation or oracle transformation is used for this enlargement.}\label{fig:experimental}\end{figure}

The Adrien experimental pair provides masked-image integration evidence. Without independent truth, accepted counts and visual remapping cannot establish displacement accuracy. We therefore do not combine those counts with synthetic correct-coverage statistics. Definitive experimental qualification requires acquisition calibration, documented speckle statistics, and independent displacement or strain evidence.

\subsection{Performance evidence}
Component time and complete scripted time are different quantities. On one synthetic native-resolution pair, an isolated paired comparison reduced median script wall time from 16.451 to 11.541 seconds by caching solver pixels. Three measured repetitions followed one warm-up and alternated arm order. Both arms obtained 1,060 correct of 1,121 visible windows, with identical states, status and provenance. These times include image generation, process startup, diagnostics and writing results; they are not kernel latency or camera throughput. The GPU was not exclusively reserved, although no competing benchmark was observed at the recorded checks.

A four-pair independent check of shared prefit and deferred rotation generation preserved 5,011 correct of 5,242 visible windows, identical acceptance masks and zero visible false acceptances. No additional rotation hypotheses were requested on these pairs. This establishes parity for that sample, not the cost of rotational fallback. Current measurements do not support a subsecond end-to-end claim. Future timing must include transfers, failed hypotheses and required result recovery, with the same images and window budgets.



\section{Discussion}
Propagation changes what is required of an initializer. One useful regional seed can enable many local solves, while a high mean matching score with an unseeded region remains insufficient. The causal ablation shows why one-pass counts alone can overstate the advantage of a costly proposal stage: the remaining affine-plus-correlation gain over affine alone is 736 centres after propagation. The confirmatory workflow adds a distinct consistency policy; its results cannot be attributed to the median/three-attempt policy used for that ablation.

The inspected model is still learned in the no-residual baseline. Its independent encoder supplies global correspondences before affine fitting, while the tested post-rectification residual checkpoint is omitted. The negative residual ablation concerns that checkpoint in a role for which it was not specifically trained. It neither rejects all learned fields nor demonstrates that a neural field is necessary. The main demonstrated benefit is the interaction of image-derived initialization and subsequent native-pixel solves.

The comparison with prior art supports a reproducible implementation and qualification contribution, not a claim to have invented affine propagation. The literature comparison in the supplement separates full-text evidence from abstracts and leaves undocumented scheduling or mask rules unspecified. Further novelty claims require full-text comparison and matched competing implementations, not a list of recent descriptors.

The method does not solve all failure modes. Extreme shear can invalidate a finite affine support even with analytical centre parameters. Smaller support changes spatial resolution and gradient uncertainty. Endpoint-mask adjacency is not a guarantee against crossing unknown interfaces. A large accepted fraction and low absolute intensity residual do not establish metrological accuracy; the International Digital Image Correlation Society guides provide the appropriate experimental reporting context \citep{guide,guide2}. The second edition's availability has been checked; no new compliance claim is inferred from that fact.

WebGPU portability concerns access to a common compute interface, not identical performance on every computer. Shared kernels reduce maintenance divergence, but host orchestration, memory ownership and solver versions still matter. The available parity checks establish preservation on specified configurations. They do not demonstrate an entirely resident pipeline or real-time operation.

\section{Conclusions}
The retained approach uses matching to propose, the photometric solver to estimate, and independent truth to evaluate. On the controlled heterogeneous ablation, status-seeded propagation raises affine-plus-correlation recovery from 16,266 to 20,651 of 22,416 visible centres. A separate frozen confirmation-and-consistency workflow recovers 23,570 of 24,641 heterogeneous centres. The policies and populations remain explicitly distinct. Severe shear, interface connectivity and gradient uncertainty bound these results. Experimental remapping illustrates operation without establishing accuracy. The work supplies an initialization route for an existing DIC solver, with portable kernel implementations and auditable limitations, rather than a replacement measurement method.

\appendix
\section{Pixel centres and composition}\label{sec:coordinates}
For native-to-reduced scale $S=\operatorname{diag}(s_x,s_y)$, the common-image coordinate map is $\mathbf z=S^{-1}(\mathbf x+\tfrac12\mathbf1)-\tfrac12\mathbf1$. It is not $\mathbf x/S$. A feature grid with $N_x\times N_y$ nodes in a $W\times H$ raster uses
\begin{equation}
z_{ij,x}=(i+\tfrac12)W/N_x-\tfrac12,\qquad
z_{ij,y}=(j+\tfrac12)H/N_y-\tfrac12.
\end{equation}
Thus $192/48=4$ places horizontal nodes at $4i+1.5$. Bilinear evaluation and spline fitting must use these same physical nodes. For equal source/target scale,
\begin{equation}
\mathbf u_n=S\mathbf u_r,\qquad F_n=SF_rS^{-1},\qquad G_n=F_n-I.
\end{equation}
The diagonal gradients are unchanged, while cross derivatives scale by $s_x/s_y$ and $s_y/s_x$. The isotropic factor 20 leaves the entire gradient invariant. The implemented bilinear evaluator clamps indices outside the nodal extent; exact affine reproduction is therefore asserted only inside that extent, not for extrapolated boundary derivatives.

More generally let $\mathbf z_0=B_0\mathbf x+\mathbf b_0$ and $\mathbf z_1=B_1\mathbf y+\mathbf b_1$ encode resizing, cropping or padding. If the reduced map is $\phi_r$, then
\begin{equation}
\phi_n(\mathbf x)=B_1^{-1}[\phi_r(B_0\mathbf x+\mathbf b_0)-\mathbf b_1],\qquad
F_n=B_1^{-1}F_r B_0.
\end{equation}
These formulas state the required conversion; they are not a claim that the tested matcher accepts differently sized frames. The present pair model requires equal shapes, with area resizing and no learned-input crop/padding in the cited synthetic configuration. Rotation and affine pullbacks are separate bilinear resamplings with explicit fill. The experimental dimensions and model raster are recorded separately in its campaign configuration.

For affine preconditioning $H(\mathbf z)=A\mathbf z+\mathbf t$ and residual map $R(\mathbf z)=\mathbf z+\mathbf r(\mathbf z)$, the forward composition is $H\circ R$ and its derivative is $A(I+\nabla\mathbf r)$. Inverse composition inside IC-GN instead uses the inverse incremental local warp; these operations must not be conflated. Analytical conversion checks in the revision package cover anisotropic factors, offsets and composition, without changing archived solver outputs.

\section*{Data and code availability}
The project is hosted at \url{https://github.com/jeffwitz/WGPU-Matching}. The inspected working tree is based on commit \texttt{4dce5d0} and contains subsequent uncommitted changes. The revision package includes source snapshots, Secure Hash Algorithm 256-bit (SHA-256) fingerprints, exact recorded campaign configurations, per-window recount scripts and reproducible figures. Missing historical solver fingerprints, external image access and a persistent release remain documented in the revision audit. The public commit alone is insufficient to reproduce all evidence.
\section*{Declarations}
Authors, affiliations, correspondence details, contribution roles, funding and competing interests require author confirmation. Acquisition provenance and permission to publish the experimental images must be confirmed. Assisted drafting and coding require author review and disclosure under publisher policy. No declaration of absence or permission is inferred.
\begin{thebibliography}{99}
\bibitem{zhou} Zhou Y, Pan B, Chen YQ. Large deformation measurement using digital image correlation: a fully automated approach. Applied Optics. 2012;51:7674--7683. doi:10.1364/AO.51.007674.
\bibitem{pan} Pan B. Reliability-guided digital image correlation for image deformation measurement. Applied Optics. 2009;48:1535--1542. doi:10.1364/AO.48.001535.
\bibitem{li} Li J, Li X, Liang Z. Enhanced feature-based path-independent initial value estimation for robust point-wise digital image correlation. Optics and Lasers in Engineering. 2019;121:189--202. doi:10.1016/j.optlaseng.2019.04.016.
\bibitem{lowe} Lowe DG. Distinctive image features from scale-invariant keypoints. International Journal of Computer Vision. 2004;60:91--110. doi:10.1023/B:VISI.0000029664.99615.94.
\bibitem{hu} Hu Y, Song R, Li Y. Efficient coarse-to-fine PatchMatch for large displacement optical flow. Proceedings of the Institute of Electrical and Electronics Engineers Conference on Computer Vision and Pattern Recognition. 2016:5704--5712. doi:10.1109/CVPR.2016.615.
\bibitem{deepmatching} Revaud J, Weinzaepfel P, Harchaoui Z, Schmid C. DeepMatching: hierarchical deformable dense matching. arXiv:1506.07656. 2015. \url{https://arxiv.org/abs/1506.07656}.
\bibitem{eloftr} Wang Y, He X, Peng S, Tan D, Zhou X. Efficient LoFTR: semi-dense local feature matching with sparse-like speed. Proceedings of the Institute of Electrical and Electronics Engineers / Computer Vision Foundation Conference on Computer Vision and Pattern Recognition. 2024:21666--21675. \url{https://openaccess.thecvf.com/content/CVPR2024/html/Wang_Efficient_LoFTR_Semi-Dense_Local_Feature_Matching_with_Sparse-Like_Speed_CVPR_2024_paper.html}.
\bibitem{affnet} Mishkin D, Radenovic F, Matas J. Repeatability is not enough: learning affine regions via discriminability. European Conference on Computer Vision. 2018. \url{https://openaccess.thecvf.com/content_ECCV_2018/html/Dmytro_Mishkin_Repeatability_Is_Not_ECCV_2018_paper.html}.
\bibitem{webgpu} World Wide Web Consortium. WebGPU specification. \url{https://www.w3.org/TR/webgpu/}. Accessed 14 November 2026.
\bibitem{wgsl} World Wide Web Consortium. WebGPU Shading Language. \url{https://www.w3.org/TR/WGSL/}. Accessed 14 November 2026.
\bibitem{guide} International Digital Image Correlation Society. Jones EMC, Iadicola MA, editors. A Good Practices Guide for Digital Image Correlation. 2018. doi:10.32720/idics/gpg.ed1.
\bibitem{dad} Edstedt J, B\"okman G, Wadenb\"ack M, Felsberg M. DaD: Distilled Reinforcement Learning for Diverse Keypoint Detection. 2025. doi:10.48550/arXiv.2503.07347.
\bibitem{dedode} Edstedt J, B\"okman G, Zhao Z. DeDoDe v2: Analyzing and Improving the DeDoDe Keypoint Detector. 2024. doi:10.48550/arXiv.2404.08928.
\bibitem{lightglue} Lindenberger P, Sarlin PE, Pollefeys M. LightGlue: Local Feature Matching at Light Speed. Proceedings of the Institute of Electrical and Electronics Engineers / Computer Vision Foundation International Conference on Computer Vision. 2023:17627--17638. doi:10.48550/arXiv.2306.13643.
\bibitem{baker} Baker S, Matthews I. Lucas--Kanade 20 years on: a unifying framework. International Journal of Computer Vision. 2004;56:221--255. doi:10.1023/B:VISI.0000011205.11775.FD.
\bibitem{panfast} Pan B, Li K, Tong W. Fast, robust and accurate digital image correlation calculation without redundant computations. Experimental Mechanics. 2013. Author manuscript: \url{https://ncorr.com/download/publications/panfast.pdf}.
\bibitem{daisy} Tola E, Lepetit V, Fua P. DAISY: an efficient dense descriptor applied to wide-baseline stereo. IEEE Transactions on Pattern Analysis and Machine Intelligence. 2010;32:815--830. doi:10.1109/TPAMI.2009.77.
\bibitem{barnes} Barnes C, Shechtman E, Finkelstein A, Goldman DB. PatchMatch: a randomized correspondence algorithm for structural image editing. ACM Transactions on Graphics. 2009;28(3). Author manuscript: \url{https://gfx.cs.princeton.edu/pubs/Barnes_2009_PAR/patchmatch.pdf}.
\bibitem{hardnet} Mishchuk A, Mishkin D, Radenovic F, Matas J. Working hard to know your neighbor's margins: local descriptor learning loss. Advances in Neural Information Processing Systems. 2017. doi:10.48550/arXiv.1705.10872.
\bibitem{ransac} Fischler MA, Bolles RC. Random sample consensus: a paradigm for model fitting with applications to image analysis and automated cartography. Communications of the ACM. 1981;24:381--395. doi:10.1145/358669.358692.
\bibitem{xfeat} Potje G, Cadar F, Araujo A, Martins R, Nascimento ER. XFeat: Accelerated Features for Lightweight Image Matching. Conference on Computer Vision and Pattern Recognition. 2024. \url{https://openaccess.thecvf.com/content/CVPR2024/html/Potje_XFeat_Accelerated_Features_for_Lightweight_Image_Matching_CVPR_2024_paper.html}.
\bibitem{guide2} International Digital Image Correlation Society. Jones EMC, Iadicola MA, editors. A Good Practices Guide for Digital Image Correlation, Edition 2. 2025. doi:10.32720/idics/gpg.ed2.
\end{thebibliography}

\end{document}
